\documentclass[10pt]{article}
\usepackage[T1]{fontenc}
\usepackage[utf8]{inputenc}
\usepackage[french]{babel}
\usepackage{amsmath,amssymb,mathrsfs,mathpazo}
\usepackage{geometry}
\geometry{a4paper,margin=2cm}
\usepackage{graphicx,array,tabularx,tabularray,colortbl,multirow,multicol,booktabs,enumitem,pifont,ulem,fancybox,diagbox,cancel}
\usepackage{pstricks,pstricks-add,pst-plot,pst-tree,pst-node,pst-text,pst-func,pst-math,pst-bspline,pst-3dplot,pst-eucl,pst-solides3d}
\let\pscalculate\undefined
\usepackage{pst-bezier}
\usepackage{tikz,pgfplots,tkz-tab}
\usetikzlibrary{arrows,arrows.meta,patterns,calc,positioning,trees,decorations.markings,shapes.geometric,shapes.misc,shapes.arrows,intersections,decorations.pathmorphing,decorations.pathreplacing,backgrounds,shadows,babel}
\pgfplotsset{compat=1.18}
\usepgfplotslibrary{fillbetween}
\usepackage[np]{numprint}
\usepackage{listings,makecell,pas-tableur,twoopt}
\let\barre\undefined
\usepackage[pstricks]{tablvar}
\definecolor{darkgreen}{RGB}{0,100,0}
\providecommand{\red}{\color{red}}
\providecommand{\blue}{\color{blue}}
\providecommand{\green}{\color{green}}
\providecommand{\black}{\color{black}}
\providecommand{\white}{\color{white}}
\providecommand{\gray}{\color{gray}}
\providecommand{\cyan}{\color{cyan}}
\providecommand{\magenta}{\color{magenta}}
\providecommand{\yellow}{\color{yellow}}
\providecommand{\violet}{\color{violet}}
\providecommand{\orange}{\color{orange}}
\providecommand{\euro}{\texteuro}
\providecommand{\EUR}{\texteuro}
\providecommand{\ieme}{\textsuperscript{e}}
\providecommand{\degres}{\ensuremath{^\circ}}
\providecommand{\degree}{\ensuremath{^\circ}}
\providecommand{\vv}[1]{\overrightarrow{#1}}
\providecommand{\Coord}[1]{#1}
\providecommand{\decosix}{}
\providecommand{\starredbullet}{\ensuremath{\bullet}}
\pagestyle{empty}
\setlength{\parindent}{0pt}
\usepackage{adjustbox}\DeclareUnicodeCharacter{2236}{:}\DeclareUnicodeCharacter{2190}{\ensuremath{\leftarrow}}
\providecommand{\R}{\mathbb{R}}
\providecommand{\N}{\mathbb{N}}
\providecommand{\Z}{\mathbb{Z}}
\providecommand{\Q}{\mathbb{Q}}
\providecommand{\D}{\mathbb{D}}
\providecommand{\C}{\mathbb{C}}
\providecommand{\e}{{\mathrm{e}}}
\providecommand{\vect}[1]{\overrightarrow{#1}}
\providecommand{\vectt}[1]{\overrightarrow{\mathrm{#1}}}
\providecommand{\vvt}[1]{\overrightarrow{\mathrm{#1}}}
\providecommand{\np}[2][]{{#2~#1}}
\providecommand{\vv}[1]{{\overrightarrow{#1}}}
\renewcommand{\d}{{\mathrm{d}}}
\providecommand{\degres}{{^\circ}}
\providecommand{\degre}{{^\circ}}
\providecommand{\blue}{\color{blue}}
\providecommand{\red}{\color{red}}
\providecommand{\enskip}{\quad}
\providecommand{\llbracket}{\lbrack\!\lbrack}
\providecommand{\rrbracket}{\rbrack\!\rbrack}
\providecommand{\euro}{{\text{€}}}
\providecommand{\no}{{\text{n°}}}
\providecommand{\interval}[2]{{[#1;#2]}}
\providecommand{\intervFO}[2]{{[#1;#2[}}
\providecommand{\dotfill}{\ldots}

\providecommand{\ds}{\displaystyle}
\providecommand{\textding}[1]{#1}
\providecommand{\Oij}{\(\mathrm{O};\vec{\imath},\vec{\jmath}\)}
\providecommand{\Oijk}{\(\mathrm{O};\vec{\imath},\vec{\jmath},\vec{k}\)}
\providecommand{\Ouv}{\(\mathrm{O};\vec{u},\vec{v}\)}
\providecommand{\intervFF}[2]{\ensuremath{[#1;#2]}}
\renewcommand{\euro}{\ifmmode\text{\texteuro}\else\texteuro\fi}
\DeclareUnicodeCharacter{00B0}{\ifmmode^{\circ}\else\textdegree\fi}
\DeclareUnicodeCharacter{20AC}{\ifmmode\text{\texteuro}\else\texteuro\fi}
% Énoncés : Éducation nationale. Transcription et corrigés : APMEP, Denis Vergès et contributeurs. Découpage et classement : CoopMaths, Sébastien Lozano, Éric Elter et contributeurs. Adaptation de la présentation : Maths Explorer.
% CC BY-SA : https://coopmaths.fr/www/about/
% https://coopmaths.fr/alea/static/bac/2024/tex/bac_2024_05_sujet2_ameriquenord_2.tex
\begin{document}
\noindent\textit{APMEP / CoopMaths. CC BY-SA. Présentation adaptée par Maths Explorer.}\par\medskip
On considère le pavé droit ABCDEFGH tel que AB $= 3$ et AD = AE $= 1$ représenté ci-dessous.

\begin{center}
\psset{unit=1cm}
\begin{pspicture}(9.4,6.3)
\psframe(0.3,0.6)(7.8,4.3)
\psline(7.8,0.6)(8.6,2.2)(8.6,5.9)(7.8,4.3)
\psline(8.6,5.9)(1.1,5.9)(0.3,4.3)
\psline[linestyle=dashed](0.3,0.6)(1.1,2.2)(8.6,2.2)
\psline[linestyle=dashed](1.1,2.2)(1.1,5.9)
\psdots[dotstyle=+,dotangle=45,dotscale=1.9](2.8,0.6)(4.85,2.2)\uput[d](2.8,0.6){I}
\uput[dl](0.3,0.6){A} \uput[dr](7.8,0.6){B} \uput[r](8.6,2.2){C} \uput[ul](1.1,2.2){D}
\uput[l](0.3,4.3){E} \uput[r](7.8,4.3){F} \uput[ur](8.6,5.9){G} \uput[ul](1.1,5.9){H}
\uput[u](4.85,2.2){M}
\end{pspicture}
\end{center}

On considère le point I du segment [AB] tel que $\vect{\text{AB}}= 3 \vect{\text{AI}}$ et on appelle $M$ le milieu du segment [CD].

On se place dans le repère orthonormé $\left(\text{A}~;~ \vect{\text{AI}}, \vect{\text{AD}}, \vect{\text{AE}}\right)$.

\medskip

\begin{enumerate}
\item Sans justifier, donner les coordonnées des points F{}, H et M.

\item 
	\begin{enumerate}
		\item 	Montrer que le vecteur $\vect{n}\begin{pmatrix}2\\6\\ 3\end{pmatrix}$ est un vecteur normal au plan (HMF).
		\item En déduire qu'une équation cartésienne du plan (HMF) est :

\[2x + 6y + 3z - 9 = 0.\]

		\item Le plan $\mathcal{P}$ dont une équation cartésienne est $5x + 15y - 3z + 7 = 0$ est-il parallèle au plan (HMF) ? Justifier la réponse.
	\end{enumerate}
\item Déterminer une représentation paramétrique de la droite (DG).
\item On appelle $N$ le point d'intersection de la droite (DG) avec le plan (HMF).

Déterminer les coordonnées du point N.
\item Le point R de coordonnées $\left(3~;~\dfrac{1}{4}~;~\dfrac{1}{2}\right)$ est-il le projeté orthogonal du point G sur le plan (HMF) ? Justifier la réponse.
\end{enumerate}

\end{document}
